\section{Relation to constructive and lattice-gauge work}\label{sec:related}

\paragraph{Heat-kernel and Villain interactions.}
Menotti and Onofri introduced the group heat-kernel action as a lattice gauge
interaction and emphasized its weak-coupling Gaussian structure
\cite{MenottiOnofri1981}.  Chevyrev and Garban recently characterized Villain
interactions in lattice gauge theory through limits of Wilson and Manton
interactions on enlarged graphs \cite{ChevyrevGarban2025}.  Our use of the
$\SUtwo$ heat kernel is finite-block and analytic: the paired Poisson images
provide explicit local derivative bounds.  We do not prove an equivalence of
regulators or a scaling limit.

\paragraph{Gauge reduction.}
Maximal-tree gauges are classical finite-lattice coordinates; they have also
been used to identify independent Hamiltonian lattice-gauge variables
\cite{LigterinkWaletBishop2000}.  Our tree gauge is not a novelty claim.  Its
role is to make the measure accounting in a computer-checkable local theorem
exact.  The based quotient leaves global conjugation visible and avoids any
claim of a global Coulomb slice.

\paragraph{Constructive renormalization.}
Balaban's program goes far beyond one block.  It develops propagators and
renormalization transformations \cite{Balaban1984}, localized small-field
effective actions and coupling renormalization in four dimensions
\cite{Balaban1987}, cluster expansions \cite{Balaban1988}, and large-field
operations used in ultraviolet stability \cite{Balaban1989}.  The present
theorem is weaker in physical reach: it neither produces a localized polymer
expansion nor controls a sequence of effective actions.  Its distinct purpose
is to expose one block as a completely explicit proof object, including exact
topology, a certified spectral floor, nonlinear coordinate bounds, and the
Hessian effect of one marginalization.

\paragraph{Convex marginalization.}
The covariance step is an application of the Brascamp--Lieb inequality
\cite{BrascampLieb1976}; the conditional-law comparison uses the implication
from logarithmic Sobolev to quadratic transport developed by Otto and Villani
\cite{OttoVillani2000}.  The formulas are standard.  What is certified here is
that the exact $\SUtwo$ block lies within their hypotheses on one stated
cylinder with rational margins.

\paragraph{Constructive gauge theory and the continuum target.}
Foundational constructive work established important inequalities,
positivity properties, and lower-dimensional or gauge--matter results
\cite{BrydgesFrohlichSeiler1979}.  These results also illustrate why a finite
local coercivity statement is not a continuum mass gap.  The theorem here
does not address the existence and spectral requirements in the Yang--Mills
Millennium problem \cite{JaffeWitten2000}.

Our literature search located no primary source stating the same package of
block, selected retained map, rational floor, local constants, and executable
certificate.  This supports a conservative claim of integrated explicitness;
it does not support a claim that the underlying methods are new or that no
similar finite computation exists.
